% GENERATED FILE. Edit canonical lessons, then run npm run generate:books.
% canonical-source-hash: fnv1a64:ed95bedfec4bfeac
% canonical-lessons: ML-C536-raktaparisodhana, ML-C536-narampu, ML-C536-pesi, ML-C536-svasakosam, ML-C536-visramikkuka

\chapter{Blood Tests and Muscles}
\label{ch:ml-blood-tests-and-muscles}
\hlchaptermodality{536}

\begin{chapteropening}
\textbf{By the end of this chapter:} \emph{I can say a blood test, a vein, a muscle, the lungs and to take a rest.}

Complete the last lesson of chapter 536: I can say a blood test, a vein, a muscle, the lungs and to take a rest.
\end{chapteropening}

\section[\texorpdfstring{\ml{രക്തപരിശോധന} (raktapariśōdhana)}{രക്തപരിശോധന (raktapariśōdhana)}]{\ml{രക്തപരിശോധന} (raktapariśōdhana) — a blood test}

\label{lesson:ML-C536-raktaparisodhana}



\begin{quote}
Before the new one: say the Malayalam for rheumatism, then the Malayalam for to throw up.
\end{quote}

\subsection*{You'll want to know: \ml{രക്തപരിശോധന}}
\textbf{\ml{രക്തപരിശോധന}} — \emph{raktapariśōdhana} — \textquotedblleft{}a blood test\textquotedblright{}.

\begin{quote}
\ml{രാവിലെ} \ml{ഒന്നും} \ml{കഴിക്കാതെ} \ml{രക്തപരിശോധന} \ml{നടത്തണം}.
\end{quote}

\begin{quote}
\emph{rāvile onnuṁ kaḻikkāte raktapariśōdhana naṭattaṇaṁ.}
\end{quote}

\begin{quote}
The blood test has to be done in the morning before eating.
\end{quote}

\ml{രക്തം} blood + \ml{പരിശോധന} a check; the verb \ml{പരിശോധിക്കുക} is already familiar.

\begin{cousinweb}
\textbf{\ml{രക്തപരിശോധന}} (\emph{raktapariśōdhana}) is a learned compound built from Sanskrit \emph{rakta}, \textquotedblleft{}red; blood\textquotedblright{} (CDIAL 10539), and \emph{pariśodhana}, \textquotedblleft{}cleaning, purifying\textquotedblright{} (Monier-Williams), which Malayalam uses for a check.
\end{cousinweb}

\begin{grammarlens}[title={What you've built}]
One more word for describing what you lived and felt.
\end{grammarlens}

\subsection*{Your turn}
\begin{itemize}
\raggedright
  \item \emph{Say it:} \emph{raktapariśōdhana}
  \item \emph{Say it:} \emph{raktapariśōdhana}, once more
  \item \emph{Say it:} \emph{rāvile onnuṁ kaḻikkāte raktapariśōdhana naṭattaṇaṁ.} — The blood test has to be done in the morning before eating.
  \item \emph{Recall:} say the Malayalam for rheumatism, then the Malayalam for to throw up, then say \emph{raktapariśōdhana} again
\end{itemize}

\begin{tcolorbox}[breakable,colback=teal!4,colframe=teal!35!black,title={Before you move on}]
What does \ml{രക്തപരിശോധന} mean? (\textquotedblleft{}A blood test\textquotedblright{}.) Say it once more.
\end{tcolorbox}
\section[\texorpdfstring{\ml{ഞരമ്പ്} (ñarampŭ)}{ഞരമ്പ് (ñarampŭ)}]{\ml{ഞരമ്പ്} (ñarampŭ) — a vein; a nerve}

\label{lesson:ML-C536-narampu}



\begin{quote}
Before the new one: say the Malayalam for to throw up, then the Malayalam for a blood test.
\end{quote}

\subsection*{You'll want to know: \ml{ഞരമ്പ്}}
\textbf{\ml{ഞരമ്പ്}} — \emph{ñarampŭ} — \textquotedblleft{}a vein; a nerve\textquotedblright{}.

\begin{quote}
\ml{നഴ്സ്} \ml{കയ്യിലെ} \ml{ഞരമ്പിൽ} \ml{നിന്ന്} \ml{രക്തം} \ml{എടുത്തു}.
\end{quote}

\begin{quote}
\emph{naḻsŭ kayyile ñarampil ninnŭ raktaṁ eṭuttu.}
\end{quote}

\begin{quote}
The nurse took blood from a vein in my arm.
\end{quote}

One word for a vein, a nerve and a sinew.

\begin{cousinweb}
\textbf{\ml{ഞരമ്പ്}} (\emph{ñarampŭ}) is native Dravidian: DEDR 2903 gives it as a sinew, a tendon, a nerve and the pulse (DEDR 2903).
\end{cousinweb}

\begin{grammarlens}[title={What you've built}]
One more word for describing what you lived and felt.
\end{grammarlens}

\subsection*{Your turn}
\begin{itemize}
\raggedright
  \item \emph{Say it:} \emph{ñarampŭ}
  \item \emph{Say it:} \emph{ñarampŭ}, once more
  \item \emph{Say it:} \emph{naḻsŭ kayyile ñarampil ninnŭ raktaṁ eṭuttu.} — The nurse took blood from a vein in my arm.
  \item \emph{Recall:} say the Malayalam for to throw up, then the Malayalam for a blood test, then say \emph{ñarampŭ} again
\end{itemize}

\begin{tcolorbox}[breakable,colback=teal!4,colframe=teal!35!black,title={Before you move on}]
What does \ml{ഞരമ്പ്} mean? (\textquotedblleft{}A vein; a nerve\textquotedblright{}.) Say it once more.
\end{tcolorbox}
\section[\texorpdfstring{\ml{പേശി} (pēśi)}{പേശി (pēśi)}]{\ml{പേശി} (pēśi) — a muscle}

\label{lesson:ML-C536-pesi}



\begin{quote}
Before the new one: say the Malayalam for a blood test, then the Malayalam for a vein.
\end{quote}

\subsection*{You'll want to know: \ml{പേശി}}
\textbf{\ml{പേശി}} — \emph{pēśi} — \textquotedblleft{}a muscle\textquotedblright{}.

\begin{quote}
\ml{വ്യായാമം} \ml{പേശികൾക്ക്} \ml{ശക്തി} \ml{കൊടുക്കും}.
\end{quote}

\begin{quote}
\emph{vyāyāmaṁ pēśikaḷkkŭ śakti koṭukkuṁ.}
\end{quote}

\begin{quote}
Exercise gives the muscles strength.
\end{quote}

Plural: \ml{പേശികൾ}, \emph{pēśikaḷ}.

Plural: \ml{പേശികൾ}.

\begin{cousinweb}
\textbf{\ml{പേശി}} (\emph{pēśi}) is borrowed from Sanskrit \emph{peśī}, \textquotedblleft{}a piece of flesh\textquotedblright{} (Monier-Williams).
\end{cousinweb}

\begin{grammarlens}[title={What you've built}]
One more word for describing what you lived and felt.
\end{grammarlens}

\subsection*{Your turn}
\begin{itemize}
\raggedright
  \item \emph{Say it:} \emph{pēśi}
  \item \emph{Say it:} \emph{pēśi}, once more
  \item \emph{Say it:} \emph{vyāyāmaṁ pēśikaḷkkŭ śakti koṭukkuṁ.} — Exercise gives the muscles strength.
  \item \emph{Recall:} say the Malayalam for a blood test, then the Malayalam for a vein, then say \emph{pēśi} again
\end{itemize}

\begin{tcolorbox}[breakable,colback=teal!4,colframe=teal!35!black,title={Before you move on}]
What does \ml{പേശി} mean? (\textquotedblleft{}A muscle\textquotedblright{}.) Say it once more.
\end{tcolorbox}
\section[\texorpdfstring{\ml{ശ്വാസകോശം} (śvāsakōśaṁ)}{ശ്വാസകോശം (śvāsakōśaṁ)}]{\ml{ശ്വാസകോശം} (śvāsakōśaṁ) — the lungs; a lung}

\label{lesson:ML-C536-svasakosam}



\begin{quote}
Before the new one: say the Malayalam for a vein, then the Malayalam for a muscle.
\end{quote}

\subsection*{You'll want to know: \ml{ശ്വാസകോശം}}
\textbf{\ml{ശ്വാസകോശം}} — \emph{śvāsakōśaṁ} — \textquotedblleft{}the lungs; a lung\textquotedblright{}.

\begin{quote}
\ml{പുകവലി} \ml{ശ്വാസകോശത്തിന്} \ml{ദോഷമാണ്}.
\end{quote}

\begin{quote}
\emph{pukavali śvāsakōśattinŭ dōṣamāṇŭ.}
\end{quote}

\begin{quote}
Smoking is bad for the lungs.
\end{quote}

\ml{ശ്വാസം} breath + \ml{കോശം} a vessel, a store.

\begin{cousinweb}
\textbf{\ml{ശ്വാസകോശം}} (\emph{śvāsakōśaṁ}) is a learned compound built from Sanskrit \emph{śvāsa}, \textquotedblleft{}breath\textquotedblright{} (CDIAL 12769), and \emph{kośa}, \textquotedblleft{}a cask, a vessel\textquotedblright{} (Monier-Williams): the store of breath.
\end{cousinweb}

\begin{grammarlens}[title={What you've built}]
One more word for describing what you lived and felt.
\end{grammarlens}

\subsection*{Your turn}
\begin{itemize}
\raggedright
  \item \emph{Say it:} \emph{śvāsakōśaṁ}
  \item \emph{Say it:} \emph{śvāsakōśaṁ}, once more
  \item \emph{Say it:} \emph{pukavali śvāsakōśattinŭ dōṣamāṇŭ.} — Smoking is bad for the lungs.
  \item \emph{Recall:} say the Malayalam for a vein, then the Malayalam for a muscle, then say \emph{śvāsakōśaṁ} again
\end{itemize}

\begin{tcolorbox}[breakable,colback=teal!4,colframe=teal!35!black,title={Before you move on}]
What does \ml{ശ്വാസകോശം} mean? (\textquotedblleft{}The lungs; a lung\textquotedblright{}.) Say it once more.
\end{tcolorbox}
\section[\texorpdfstring{\ml{വിശ്രമിക്കുക} (viśramikkuka)}{വിശ്രമിക്കുക (viśramikkuka)}]{\ml{വിശ്രമിക്കുക} (viśramikkuka) — to take a rest; to rest}

\label{lesson:ML-C536-visramikkuka}



\begin{quote}
Before the new one: say the Malayalam for a muscle, then the Malayalam for the lungs.
\end{quote}

\subsection*{You'll want to know: \ml{വിശ്രമിക്കുക}}
\textbf{\ml{വിശ്രമിക്കുക}} — \emph{viśramikkuka} — \textquotedblleft{}to take a rest; to rest\textquotedblright{}.

\begin{quote}
\ml{ഉച്ചയ്ക്ക്} \ml{ശേഷം} \ml{മുത്തച്ഛൻ} \ml{കുറച്ച്} \ml{നേരം} \ml{വിശ്രമിക്കും}.
\end{quote}

\begin{quote}
\emph{uccaykkŭ śēṣaṁ muttacchan kuṟaccŭ nēraṁ viśramikkuṁ.}
\end{quote}

\begin{quote}
After noon grandfather rests for a while.
\end{quote}

Past: \ml{വിശ്രമിച്ചു}, \emph{viśramiccu}.

The noun \ml{വിശ്രമം}, rest, is already familiar.

\begin{cousinweb}
\textbf{\ml{വിശ്രമിക്കുക}} (\emph{viśramikkuka}) is built from Sanskrit \emph{viśrama}, \textquotedblleft{}rest\textquotedblright{} (CDIAL 11947), with the Malayalam verb ending -\ml{ിക്കുക}.
\end{cousinweb}

\begin{grammarlens}[title={What you've built}]
One more word for describing what you lived and felt.
\end{grammarlens}

\subsection*{Your turn}
\begin{itemize}
\raggedright
  \item \emph{Say it:} \emph{viśramikkuka}
  \item \emph{Say it:} \emph{viśramikkuka}, once more
  \item \emph{Say it:} \emph{uccaykkŭ śēṣaṁ muttacchan kuṟaccŭ nēraṁ viśramikkuṁ.} — After noon grandfather rests for a while.
  \item \emph{Recall:} say the Malayalam for a muscle, then the Malayalam for the lungs, then say \emph{viśramikkuka} again
\end{itemize}

\begin{tcolorbox}[breakable,colback=teal!4,colframe=teal!35!black,title={Before you move on}]
What does \ml{വിശ്രമിക്കുക} mean? (\textquotedblleft{}To take a rest; to rest\textquotedblright{}.) Say it once more.
\end{tcolorbox}
